\chapter{Experimental Methods}
\label{ch:methods}

\section{Common conventions}

All controlled questions are assessed using a three-class decision space: original answer 1, original answer 2, or tie. A parser may also encounter an invalid or unavailable response, but invalid and terminal states are not silently treated as ties. For an estimator requiring a complete or decisive case, the denominator is reported explicitly. The project uses deterministic analysis seeds and 10,000 bootstrap iterations for confidence intervals. Unless noted otherwise, uncertainty intervals are non-parametric bootstrap 95\% intervals from the frozen analysis.

The frozen \texttt{controlled-judge-pairwise-v1} prompt instructs an impartial comparison of the question and two candidate answers and requires a JSON-only response. The required fields are \texttt{verdict}, \texttt{criteria\_scores}, \texttt{confidence}, and \texttt{explanation}. Verdict admits \texttt{ANSWER\_A}, \texttt{ANSWER\_B}, \texttt{TIE}, and \texttt{UNKNOWN}; only the first three are valid scientific labels. Correctness, relevance, completeness, clarity, and safety scores are strict integers 1--5; confidence is a finite self-reported number in [0,1]; explanation is nonempty. The adapter may remove a code fence and isolate an apparent JSON object, but strict schema validation rejects malformed, missing, extra, or out-of-range fields as \texttt{INVALID\_RESPONSE}, without heuristic relabeling. Displayed A/B verdicts are mapped through persisted presented answer IDs to canonical answer identity before comparison; ties remain ties.

The source-corrected pass manifest preserves temperatures \texttt{0.0000} or \texttt{0.7000}, top-p \texttt{1.0000}, and provider-request seeds \texttt{null}, \texttt{42}, or \texttt{20260818}; non-null request seeds occur for GPT only. RQ2 plans five repetitions per exact fixed-configuration group. The separate Multi-Judge manifest specifies temperature 0, top-p 1, maximum output tokens 350, and no provider request seed. Analysis/bootstrap seeds are 20260818 for RQ1--RQ7 Primary and RQ6, and 20260823 for Multi-Judge; neither is a universal provider decoding seed.

The frozen \texttt{controlled-retry-v2} policy permits up to two retries for rate limits, temporary HTTP 5xx, or connection failures with 2- and 8-second backoffs; timeout permits one retry after 5 seconds. Invalid responses, refusals, authentication/configuration failures, provenance mismatches, budget blocks, unsupported parameters, and ambiguous attempts are terminal. Attempts are persisted before sending. Bounded recovery preserves scientific inputs and never duplicates a successful pass. A retry is an operational attempt, not an additional RQ2 repetition. Fourteen terminal Claude observations remained unavailable and are handled by estimator-specific complete-case rules without imputation.

The study does not run a new provider call during analysis. It reads persisted observations with their state and provenance. This distinction is useful because an analysis result remains reviewable even if a provider is unavailable later. It also prevents a rerun with a changed model or routing policy from silently replacing the original evidence.

\section{RQ1: agreement with human preference reference labels}

RQ1 compares a judge outcome with the available human preference reference label on the same pair. Exact agreement is the proportion of matching three-class outcomes. Cohen's $\kappa$ is reported as a chance-corrected companion statistic. The denominator excludes observations that cannot be validly classified under the frozen policy. The result should be interpreted as agreement with a reference labeling process, not as accuracy against ground truth.

If $J_i$ is the valid judge outcome and $H_i$ the human reference label for observation $i$, exact agreement is
\[
\hat{A}=\frac{1}{n}\sum_{i=1}^{n}\mathbb{I}(J_i=H_i).
\]
The final RQ1 analytical population is 772 observations. It is smaller than the planned 800 because the analysis preserves validity and complete-case requirements rather than filling unavailable outcomes. Category and judge stratifications are descriptive extensions, not alternative primary estimands.

\section{RQ2: fixed-condition repetition consistency}

RQ2 repeats the same evaluation under the same frozen condition. It does not estimate the effect of arbitrary sampling temperatures or provider changes. The strict complete estimator asks whether all required repetitions in a complete cell agree exactly. The final result uses 1,475 strict valid cells with 7,375 valid repeated decisions. A conditional sensitivity denominator is separately retained so that readers can see how requiring complete cells changes the population.

Repeated consistency is operationally useful: it measures whether a protocol returns the same outcome when asked again under the same conditions. It should not be interpreted as an intrinsic property of a model across all prompts, all deployment settings, or all future versions. A high repetition score can coexist with substantial position sensitivity, which is exactly why RQ2 and RQ3 are not merged into a single reliability score.

\section{RQ3: answer-order sensitivity}

For each candidate pair, the study obtains an AB pass and a BA pass. The raw slot labels are converted back to the identity of the original answer. A mapped decisive flip occurs if both passes are decisive and their mapped winners differ. This mapping is essential: comparing raw first/second slot labels would confound a deliberate order swap with a preference for an original answer.

Let $M_{AB}$ and $M_{BA}$ be the mapped original-answer outcomes. The decisive flip estimator is
\[
\hat{F}=\frac{\sum_i \mathbb{I}(M_{AB,i}\neq M_{BA,i})\mathbb{I}(M_{AB,i},M_{BA,i}\in\{A,B\})}{\sum_i \mathbb{I}(M_{AB,i},M_{BA,i}\in\{A,B\})}.
\]
The analysis also retains all-paired outcome disagreement as a descriptive quantity, because ties and invalid outcomes can make the broader paired population larger than the decisive subset. The primary RQ3 result is the decisive mapped flip rate.

\section{RQ4: controlled redundant-length effect}

RQ4 uses a frozen redundant-text treatment. The treatment is intentionally narrow: it changes the defined answer variant in a controlled way while keeping the underlying pair design fixed. A stable treatment win requires the relevant matched outcomes to meet the frozen validity and stability criterion. The numerator counts only cases where the treatment variant wins; the denominator counts eligible controlled pairs.

This design answers a precise question. It does not establish that all longer answers are preferred, that length never matters, or that a judge is generally free from verbosity-related behaviour. The value of a null finding here is methodological: it bounds the claim to the implemented treatment and population.

\section{RQ5: controlled presentation-format effect}

RQ5 parallels RQ4 but uses a frozen presentation/list-prefix treatment. The same discipline applies. The outcome is a stable treatment-variant win rate on eligible controlled pairs, not a general measure of style preference. Reporting the construction, eligibility rule, and denominator is especially important for a result at or near zero because a zero numerator is meaningful only relative to the number and quality of opportunities for it to occur.

\section{RQ6: counterbalanced source-family association}

RQ6 evaluates whether the answer produced by the same source family as the judge is selected in a counterbalanced experiment. The counterbalanced lineage includes 480 planned units and uses stored exact answer/source identities. Its analytical stable decisive denominator is 159, yielding coverage of 33.13\%. The estimator is the proportion of stable decisive outcomes selecting the same-family answer. The counterbalancing design improves interpretability relative to an unbalanced display, but it does not prove a causal self-bias mechanism. Judge-level heterogeneity is reported because an aggregate near 50\% can mask very different patterns across judges.

\section{RQ7: complementary mitigation strategies}

RQ7 contains two separately defined mitigation families. The primary DUAL\_SWAP protocol is a within-judge presentation-consistency filter. It compares baseline and DUAL\_SWAP agreement on its own linked retained population and separately reports coverage relative to complete linked triplets. Because the protocol filters cases, a higher agreement alone would not be enough; the coverage change is part of the result.

The secondary Multi-Judge Consensus protocol aggregates strict four-judge votes. It retains pairs satisfying its own validity condition and compares the consensus outcome with an equal-weight individual-judge baseline on exactly the same retained pairs. Its planned population, retained population, comparator, and confidence interval are therefore distinct from DUAL\_SWAP. No direct head-to-head delta is defined between the two methods.

\begin{table}[H]
\centering
\caption{Final analysis populations by research question.}
\label{tab:populations}
\begin{tabularx}{\textwidth}{@{}p{0.11\textwidth}p{0.34\textwidth}r r X@{}}
\toprule
RQ & Estimator population & Planned & Analytical & Interpretation boundary\\
\midrule
RQ1 & valid human-label comparisons & 800 & 772 & reference-label agreement\\
RQ2 & strict complete repetition cells & 1,600 & 1,475 & fixed-condition repetition\\
RQ3 & decisive mapped AB/BA pairs & 800 & 549 & mapped order flips\\
RQ4 & eligible stable redundant-text pairs & 800 & 682 & one frozen treatment\\
RQ5 & eligible stable format pairs & 800 & 716 & one frozen treatment\\
RQ6 & stable decisive counterbalanced units & 480 & 159 & association; heterogeneity\\
RQ7 primary & retained linked DUAL\_SWAP units & 800 & 597 & within-judge filter\\
RQ7 secondary & strict consensus retained pairs & 1,611 & 1,125 & cross-judge aggregation\\
\bottomrule
\end{tabularx}
\end{table}

\section{Statistical interpretation}

Point estimates are reported with their denominators and bootstrap intervals where the final artifact provides them. The purpose of an interval is not merely to decorate a percentage. It communicates how much the observed sample constrains the estimator under the frozen resampling procedure. A confidence interval crossing zero for a matched difference is a reason to avoid a definitive positive/negative claim, even if the point estimate is positive.

Observations derived from the same original MT-Bench question may be correlated across turns, judges, repetitions, swaps, variants, and matched units. A separate deterministic 10,000-replicate question-clustered bootstrap sensitivity analysis (seed 20260912) resampled original question IDs and preserved all eligible observations belonging to each sampled question together. All substantive conclusions were unchanged. This sensitivity does not replace official frozen confidence intervals. In particular, its RQ7 Primary coverage-delta interval of -24.62 to -17.64 percentage points is sensitivity evidence only; the frozen result had no official interval for that quantity.

The project deliberately does not collapse all RQs into a single aggregate score. Agreement, repeated consistency, order sensitivity, controlled treatment response, source-family association, coverage, and consensus behaviour have different units and failure modes. Combining them would hide the trade-offs that a responsible evaluator deployment must expose.
